% Section 7.2, from lectures/L07/appendix/a_conv_layer.md (Set 1), b_max_pool_layer.md (Set 2) and
% c_flatten_layer.md (Set 3). The reference implementations in Python and C are in the repository.

\section{Exercises}
\label{c7:sec:exercises}

\begin{aside}
\textbf{How to work these.} Each set is one demo program that is already in the companion
repository, with a skeleton of the struct you write, its helper functions, the \code{main()} that
uses it, and a Makefile: \file{lectures/L06/conv_layer/cpp/conv_demo.cpp} for Set~1, which you
started in Exercise~\ref{c6:ex:3}, \file{lectures/L07/max_pool_layer/cpp/max_pool_demo.cpp} for
Set~2, and \file{lectures/L07/flatten_layer/cpp/flatten_demo.cpp} for Set~3. Until you switch the
body of \code{main()} on, as each set's last exercise describes, that body is compiled out, and with
it every use of the struct, so the program builds while the struct is still incomplete.

\textbf{How to check your work.} These three have no test suite. The conv layer's demo prints each
result next to the value Exercises~\ref{c6:ex:1} and~\ref{c6:ex:2} arrive at by hand. For the other
two, run the Python reference beside your program: it feeds the same input through the same layer
and prints the same matrices, so compare the numbers rather than the formatting.

The reference implementations are published in the repository beside each C++ demo, in Python in
\file{python} and in C in \file{c}; the C versions keep each layer behind an opaque struct whose
accessor functions are only exposed through its header. Try each exercise yourself before looking
at them.
\end{aside}

% ------------------------------------------------------------------------------------------
\exerciseset{Finishing the Conv Layer}
\label{c7:set:1}\label{c7:app:a}
You started implementing \code{ml::ConvLayer} in Exercise~\ref{c6:ex:3}: finish it here. The full
task description and the target \code{main()} usage are both in Exercise~\ref{c6:ex:3}; this set
only covers the wrap-up steps for compiling and running it for the first time, in
\file{lectures/L06/conv_layer/cpp/conv_demo.cpp}.

\codeexercise{Compiling and running the program}
\label{c7:ex:1}
As usual, the program can be run by typing the command \code{make} in the terminal:

\begin{shell}
make
\end{shell}

\noindent Once the implementation is complete, uncomment the compiler flag \code{CXX_FLAGS} in the
Makefile, \file{lectures/L06/conv_layer/cpp/Makefile}. That is, change the following:

\begin{makecode}
# TODO: Uncomment this line once the implementation is finished.
#CXX_FLAGS := -std=c++17 -Wall -Werror -DCONV_LAYER_IMPLEMENTED
\end{makecode}

\noindent to

\begin{makecode}
CXX_FLAGS := -std=c++17 -Wall -Werror -DCONV_LAYER_IMPLEMENTED
\end{makecode}

\noindent You can then also remove the conditional-compilation guard \code{CONV_LAYER_IMPLEMENTED}
from \file{conv_demo.cpp} if you like. In that case, change the following:

\begin{cppcode}
/**
 * @brief Create and demonstrate a simple convolutional layer.
 *
 * @return 0 on success, -1 on failure.
 */
int main()
{
//! @todo Remove this guard (and/or uncomment the compiler flags in the makefile) once the
//        implementation is finished.
#ifdef CONV_LAYER_IMPLEMENTED

    // Function content.
    return 0;

//! @todo Remove this guard (and/or uncomment the compiler flags in the makefile) once the
//        implementation is finished.
#endif /** CONV_LAYER_IMPLEMENTED */
}
\end{cppcode}

\noindent to

\begin{cppcode}
/**
 * @brief Create and demonstrate a simple convolutional layer.
 *
 * @return 0 on success, -1 on failure.
 */
int main()
{
    // Function content.
    return 0;
}
\end{cppcode}

% ------------------------------------------------------------------------------------------
\exerciseset{A Simple Max Pooling Layer}
\label{c7:set:2}\label{c7:app:b}

\codeexercise{The \code{ml::MaxPoolLayer} struct}
\label{c7:ex:2}
A struct named \code{ml::MaxPoolLayer} should be added to
\file{lectures/L07/max_pool_layer/cpp/max_pool_demo.cpp} to implement a simple max pooling layer.
To keep things as simple as possible, we implement a struct and skip get/set methods, deletion of
copy and move constructors, and so on.

Study the code in the \code{main()} function. Your implementation should be written so this code
works to create and use a max pooling layer named \code{poolLayer}:

\begin{cppcode}
// Create a max pooling layer: 4x4 input, 2x2 pooling regions, produces 2x2 output.
constexpr std::size_t inputSize{4U};
constexpr std::size_t poolSize{2U};
ml::MaxPoolLayer poolLayer{inputSize, poolSize};

// Example 4x4 input matrix (could represent an image or feature map).
const Matrix2d input{{2, 1, 6, 1},
                     {3, 0, 4, 6},
                     {1, 2, 4, 5},
                     {3, 4, 7, 7}};

// Perform feedforward (pooling).
poolLayer.feedforward(input);

// Example output gradients (same shape as pooling output, used for backpropagation demo).
const Matrix2d outputGradients{{1, 2},
                               {3, 4}};
// Perform backpropagation.
poolLayer.backpropagate(outputGradients);
\end{cppcode}

\codeexercise{Compiling and running the program}
\label{c7:ex:3}
As usual, the program can be run by typing the command \code{make} in the terminal:

\begin{shell}
make
\end{shell}

\noindent Once the implementation is complete, uncomment the compiler flag \code{CXX_FLAGS} in the
Makefile, \file{lectures/L07/max_pool_layer/cpp/Makefile}. That is, change the following:

\begin{makecode}
# TODO: Uncomment this line once the implementation is finished.
#CXX_FLAGS := -std=c++17 -Wall -Werror -DMAX_POOL_LAYER_IMPLEMENTED
\end{makecode}

\noindent to

\begin{makecode}
CXX_FLAGS := -std=c++17 -Wall -Werror -DMAX_POOL_LAYER_IMPLEMENTED
\end{makecode}

\noindent You can then also remove the conditional-compilation guard
\code{MAX_POOL_LAYER_IMPLEMENTED} from \file{max_pool_demo.cpp} if you like. In that case, change
the following:

\begin{cppcode}
/**
 * @brief Create and demonstrate a simple max pooling layer.
 *
 * @return 0 on success, -1 on failure.
 */
int main()
{
//! @todo Remove this guard (and/or uncomment the compiler flags in the makefile) once the
//        implementation is finished.
#ifdef MAX_POOL_LAYER_IMPLEMENTED

    // Function content.
    return 0;

//! @todo Remove this guard (and/or uncomment the compiler flags in the makefile) once the
//        implementation is finished.
#endif /** MAX_POOL_LAYER_IMPLEMENTED */
}
\end{cppcode}

\noindent to

\begin{cppcode}
/**
 * @brief Create and demonstrate a simple max pooling layer.
 *
 * @return 0 on success, -1 on failure.
 */
int main()
{
    // Function content.
    return 0;
}
\end{cppcode}

% ------------------------------------------------------------------------------------------
\exerciseset{A Simple Flatten Layer}
\label{c7:set:3}\label{c7:app:c}

\codeexercise{The \code{ml::FlattenLayer} struct}
\label{c7:ex:4}
A struct named \code{ml::FlattenLayer} should be added to
\file{lectures/L07/flatten_layer/cpp/flatten_demo.cpp} to implement a simple flatten layer. To keep
things as simple as possible, we implement a struct and skip get/set methods, deletion of copy and
move constructors, and so on.

Study the code in the \code{main()} function. Your implementation should be written so this code
works to create and use a flatten layer named \code{flattenLayer}:

\begin{cppcode}
// Create a flatten layer: 4x4 input, produces 1x16 output.
constexpr std::size_t inputSize{4U};
ml::FlattenLayer flattenLayer{inputSize};

// Example 4x4 input matrix (could represent an image or feature map).
const Matrix2d input{{2, 1, 6, 1},
                     {3, 0, 4, 6},
                     {1, 2, 4, 5},
                     {3, 4, 7, 7}};

// Perform feedforward (flatten the input).
flattenLayer.feedforward(input);

// Example output gradients (same shape as flattened output, used for backpropagation demo).
const Matrix1d outputGradients{1, 2, 3, 4, 8, 7, 6, 5, 0, 2, 4, 8, 9, 7, 5, 3};

// Perform backpropagation (unflatten the output gradients).
flattenLayer.backpropagate(outputGradients);
\end{cppcode}

\codeexercise{Compiling and running the program}
\label{c7:ex:5}
As usual, the program can be run by typing the command \code{make} in the terminal:

\begin{shell}
make
\end{shell}

\noindent Once the implementation is complete, uncomment the compiler flag \code{CXX_FLAGS} in the
Makefile, \file{lectures/L07/flatten_layer/cpp/Makefile}. That is, change the following:

\begin{makecode}
# TODO: Uncomment this line once the implementation is finished.
#CXX_FLAGS := -std=c++17 -Wall -Werror -DFLATTEN_LAYER_IMPLEMENTED
\end{makecode}

\noindent to

\begin{makecode}
CXX_FLAGS := -std=c++17 -Wall -Werror -DFLATTEN_LAYER_IMPLEMENTED
\end{makecode}

\noindent You can then also remove the conditional-compilation guard
\code{FLATTEN_LAYER_IMPLEMENTED} from \file{flatten_demo.cpp} if you like. In that case, change the
following:

\begin{cppcode}
/**
 * @brief Create and demonstrate a simple flatten layer.
 *
 * @return 0 on success, -1 on failure.
 */
int main()
{
//! @todo Remove this guard (and/or uncomment the compiler flags in the makefile) once the
//        implementation is finished.
#ifdef FLATTEN_LAYER_IMPLEMENTED
    // Function content.
    return 0;

//! @todo Remove this guard (and/or uncomment the compiler flags in the makefile) once the
//        implementation is finished.
#endif /** FLATTEN_LAYER_IMPLEMENTED */
}
\end{cppcode}

\noindent to

\begin{cppcode}
/**
 * @brief Create and demonstrate a simple flatten layer.
 *
 * @return 0 on success, -1 on failure.
 */
int main()
{
    // Function content.
    return 0;
}
\end{cppcode}
